\section{Time lower bounds from the space hierarchy}
\label{sec:hierarchy}

Theorem~\ref{thm:main} has hierarchy consequences analogous to those
Williams derives from his square-root simulation
\cite[Corollaries~1.2--1.3 and Section~4]{Williams2025}.
Write \(\mathrm{TIME}_{1w,\mathrm{ro}}(t)\) for languages decided in
worst-case \(O(t(n))\) time by the fixed deterministic machine model of
Theorem~\ref{thm:main}: one writable tape and head, unit head movements,
and a fixed number of read-only input heads, with the ordinary
cap-independent input encoding described in Section~\ref{sec:introduction}.
Write \(\mathrm{DSPACE}(s)\) for ordinary deterministic \(O(s(n))\)
work-space bits with an endmarked read-only input tape whose head stays
between the endmarkers. We use the standard
space-constructibility convention that a fixed machine can mark exactly
\(s(n)\) work cells on every input of length \(n\); constant-factor
variants give the same statements.

\begin{corollary}[One-writable-tape separation]
\label{cor:space-hierarchy}
For every space-constructible \(s(n)\ge n\) and every fixed
\(0<\epsilon<5/2\),
\[
 \mathrm{DSPACE}(s(n))\not\subseteq
 \mathrm{TIME}_{1w,\mathrm{ro}}\!\left(s(n)^{5/2-\epsilon}\right).
\]
\end{corollary}

\begin{proof}
Set \(a=5/2-\epsilon\) and
\[
 \beta=\frac25\max\{1,a\}
       =\max\left\{\frac25,1-\frac{2\epsilon}{5}\right\}<1.
\]
On an input of length \(n\), compute \(n\) with a logarithmic-space
counter and try doubling caps starting at \(T_0=\max\{2,n\}\), reusing
space after each timeout and returning only on a halt. At every trial
cap \(T\ge n\), rescanning supplies every symbol within the cap radius
in \(O(\log(n+2)+\log(T+2))=O(\log(T+2))\) work bits. Thus the access
condition holds for each capped call, even when \(a<1\).
For a run of length \(t\), the final cap is at most
\(2\max\{2,n,t\}\). Theorem~\ref{thm:main} therefore gives, for any
fixed machine with \(t(n)=O(s(n)^a)\), a uniform deterministic simulation
whose work-space usage in bits satisfies
\begin{align*}
 \text{space}
 &=O\!\left((n+t(n)+2)^{2/5}\log^C(n+t(n)+2)\right)\\
 &=O\!\left(s(n)^\beta\log^{C'}(s(n)+2)\right).
\end{align*}
This includes the length counter. No value of \(s(n)\) or running-time
bound is needed by the simulator.

Choose one \(q\) with \(\beta<q<1\). Every fixed logarithmic exponent
and constant in this bound is absorbed by \(s(n)^q\), so
\(\mathrm{TIME}_{1w,\mathrm{ro}}(s^a)\) is contained in
\(\mathrm{DSPACE}(\lceil s^q\rceil)\). The deterministic space hierarchy,
in its standard constant-factor form, separates
\(\mathrm{DSPACE}(s)\) from \(\mathrm{DSPACE}(r)\) whenever the larger
bound \(s\ge n\) is space-constructible and \(r=o(s)\); constructibility
of \(r\) is not required \cite[Theorem~3.1]{LadnerLynch1976}.
For this application, the passage from the cited exact-cell theorem
to bit-space bounds can be seen directly. Choose \(q'\) with
\(q<q'<1\). After a fixed normalization of the machine and constructor
models, an exactly constructible cell bound \(S=\Theta(s)\) is available,
and each fixed \(O(s^q)\)-bit decider uses
\(O(s^q+\log(n+2))=o(S^{q'})\) cells. For each such decider, the bound
\(\lceil S^{q'}\rceil\) therefore holds at all sufficiently large lengths;
the finitely many exceptional inputs can be handled in finite control.
Apply the exact-cell theorem with
the smaller bound \(\lceil S^{q'}\rceil=o(S)\), which need not be
constructible. Its separating machine uses \(O(S)=O(s)\) bits.
This proves the required separation from
\(\mathrm{DSPACE}(\lceil s^q\rceil)\).
\end{proof}

For a concrete example, fix a standard binary encoding of deterministic
machines \(M\) with one endmarked read-only input tape whose unit-movement
head stays between the endmarkers, one initially blank work tape over a
fixed binary-plus-blank alphabet with one unit-movement head, and an
explicit finite-control transition
table. Use a tuple encoding that stores its
fields verbatim with linear overhead and whose field symbols are
obtainable with logarithmic-space counters. Let \(H\) consist of the
encodings \(\langle M,x,1^k\rangle\)
such that \(|M|\le k\) and \(M(x)\) halts after visiting at most \(k\)
work cells; malformed codes are outside \(H\).

\begin{corollary}[Bounded-space halting]
\label{cor:bounded-space-halting}
The language \(H\) is complete for \(\mathrm{DSPACE}(n)\) under
logarithmic-space many-one reductions of linear output length.
If \(N\) denotes its input length, then for every fixed
\(0<\epsilon<5/2\),
\[
 H\notin
 \mathrm{TIME}_{1w,\mathrm{ro}}\!\left(N^{5/2-\epsilon}\right).
\]
\end{corollary}

\begin{proof}
For membership in linear space, first reject any malformed tuple or transition
table, and reject if \(|M|>k\). A universal simulation then stores the
bounded work interval, head positions and state in
\(O(k+|M|+\log(N+2))=O(N)\) bits, looking up transitions by rescanning
the code. For fixed encoding constants \(K_0,K_1\), the number of bounded
configurations is at most
\[
 K_0(N+2)2^{|M|}(k+1)^3 3^k
 \le 2^{K_1(k+|M|+\log(N+2)+1)}.
\]
The factors account for the bounded input head, state code, work-interval
and head markers, and the fixed three-symbol work alphabet. A conservative
counter for this bound uses \(O(N)\) bits. Reject an attempted visit to a
\((k+1)\)-st work cell, accept a halt within the bound, and otherwise
reject when the clock is exhausted. In the last case a deterministic
bounded configuration has repeated.

Fix \(A\in\mathrm{DSPACE}(n)\). Constant-factor space normalization
gives a fixed decider \(M_A\) in the chosen format using at most \(cn+d\)
work cells. Replace rejection by a stationary loop, leaving acceptance
as a halt, to obtain \(M_A^+\). Choose \(k=c'n+d'\) large enough to
bound this computation and satisfy \(k\ge |M_A^+|\), including small
inputs. Put \(y(x)=\langle M_A^+,x,1^k\rangle\). Then \(x\in A\) if and
only if \(y(x)\in H\), and \(|y(x)|=O(n+1)\).
This is a logarithmic-space reduction of linear output length.
More specifically, given a coordinate \(j\), counters compute the field
boundaries and return a fixed-code, header or unary symbol, or rescan
\(x\) for its indexed bit; outside the encoding they return the prescribed
blank or endmarker. This uses
\(O(\log(n+2)+\log(|j|+2))\) work bits and also computes \(|y(x)|\).
The encoded string \(y(x)\) is independent of the simulation cap.

Suppose a fixed target-model machine decided \(H\) in \(O(N^a)\) time,
where \(a=5/2-\epsilon\). Simulate it on the virtual input \(y(x)\) by
Theorem~\ref{thm:main}, supplying the symbols of its prescribed initial
writable and read-only tapes from \(y(x)\), or the fixed blank symbol,
by the preceding procedure. Start doubling caps at
\(\max\{2,|y(x)|\}\). Since \(|y(x)|\ge n\), symbol access uses
\(O(\log(T+2))\) bits at every coordinate within distance \(T\) of
the fixed head origins, for each trial cap \(T\). Thus every capped call
satisfies the access condition. The last cap is
\(O(n+1+(n+1)^a)\), giving
\[
 O\!\left((n+1)^\beta\log^C(n+2)\right)=o(n),
 \qquad \beta=\frac25\max\{1,a\}<1.
\]
Choosing one \(q\in(\beta,1)\) as above would put every
\(A\in\mathrm{DSPACE}(n)\) in \(\mathrm{DSPACE}(\lceil n^q\rceil)\),
contrary to the space hierarchy. This uses the public-symbol interface
of Theorem~\ref{thm:main}, not closure of one-tape time under
logarithmic-space reductions.
\end{proof}

These are class nonmembership statements in the one-writable-head model,
not lower bounds on every input length. They assert no general multitape
or random-access bound and no logarithmic endpoint.
